% !TEX program = xelatex
% !TEX TS-program = xelatex
% !TEX encoding = UTF-8 Unicode
% ---------------------------------------------------------------
% 本文含中文，必须使用 XeLaTeX（或 LuaLaTeX）编译。
% 用 pdfLaTeX 编译会报：CTeX fontset `mac' is unavailable in current mode
% 命令行编译：latexmk -xelatex Transformer.tex
%
% 所有插图均由 TikZ 绘制，无外部图片依赖，文末不设参考文献。
% 文中采用行向量记法；省略 batch 维度，核心推导省略 dropout 和部分偏置。
% ---------------------------------------------------------------
\RequirePackage{iftex}
\ifPDFTeX
  \errmessage{This document must be compiled with XeLaTeX (Chinese fontset is unavailable under pdfLaTeX). Please run: latexmk -xelatex Transformer.tex}
  \endinput
\fi

\documentclass[UTF8,a4paper,12pt]{ctexart}

\usepackage{amsmath}
\numberwithin{equation}{section}
\usepackage{amssymb}
\usepackage{bm}
\usepackage{graphicx}
\graphicspath{{images/}}
\usepackage{float}
\usepackage{caption}
\captionsetup{font=small,labelfont=normalfont,labelsep=colon,skip=5pt}
\usepackage{booktabs}
\usepackage{tabularx}
\usepackage{array}
\usepackage{tikz}
\usetikzlibrary{arrows.meta,positioning,calc,fit,backgrounds,shapes.misc}
\usepackage{fancyhdr}
\pagestyle{fancy}
\fancyhf{}
\fancyhead[R]{\thepage}
\renewcommand{\headrulewidth}{0pt}
\setlength{\headheight}{15pt}
\fancypagestyle{plain}{
\fancyhf{}
\fancyhead[R]{\thepage}
\renewcommand{\headrulewidth}{0pt}
}

\renewcommand{\topfraction}{0.9}
\renewcommand{\bottomfraction}{0.7}
\renewcommand{\textfraction}{0.1}
\renewcommand{\floatpagefraction}{0.8}

% 公式上下间距：必须放在 \AtBeginDocument 里，否则会被 \document
% 开头的 \normalsize 重置回类默认值（12pt），设了等于白设。
\AtBeginDocument{%
  \setlength{\abovedisplayskip}{16pt plus 3pt minus 5pt}%
  \setlength{\belowdisplayskip}{16pt plus 3pt minus 5pt}%
  \setlength{\abovedisplayshortskip}{8pt plus 2pt minus 2pt}%
  \setlength{\belowdisplayshortskip}{8pt plus 2pt minus 2pt}%
  \setlength{\jot}{6pt}%  align 等多行公式行间距（默认 3pt）
}
\setlength{\textfloatsep}{8pt plus 2pt minus 2pt}
\setlength{\intextsep}{8pt plus 2pt minus 2pt}
\setlength{\abovecaptionskip}{4pt}
% 段末安全网：仅在某段本会溢出右边距时生效，正常段落不受影响
\setlength{\emergencystretch}{1em}

\usepackage{xcolor}
\usepackage[colorlinks=true,linkcolor=black,urlcolor=blue,citecolor=black]{hyperref}

\title{Transformer}
\date{2026年9月28日}

% ---- 插图里反复用到的样式 ----
\newcolumntype{Y}{>{\raggedright\arraybackslash}X}
\DeclareMathOperator{\softmax}{softmax}
\DeclareMathOperator{\ReLU}{ReLU}
\DeclareMathOperator{\LN}{LN}
\DeclareMathOperator{\FFN}{FFN}
\DeclareMathOperator{\MHA}{MHA}
\DeclareMathOperator{\Concat}{Concat}
\DeclareMathOperator*{\argmax}{arg\,max}
\newcommand{\R}{\mathbb{R}}
\newcommand{\bos}{\texttt{<BOS>}}
\newcommand{\eos}{\texttt{<EOS>}}
\newcommand{\keypoint}[1]{\par\smallskip\noindent\textbf{#1}\par\smallskip}
\tikzset{
  flowbox/.style={rounded corners=4pt,draw=blue!55!black,fill=blue!7,
    inner sep=5pt,align=center,font=\small},
  flowarr/.style={-{Stealth[length=1.8mm,width=1.2mm]},draw=gray!75,line width=0.7pt},
  flowlab/.style={font=\small,align=center,text=black!85},
  addnorm/.style={flowbox,fill=orange!10,draw=orange!55!black},
  sumcircle/.style={circle,draw=gray!80,fill=white,inner sep=0pt,minimum size=5mm},
  tinylabel/.style={font=\scriptsize,align=center,text=black!80}
}

\begin{document}
\maketitle
\section{引言}

如今主流的大语言模型几乎都以 Transformer 为骨架，最典型的就是 GPT——全称 \textit{Generative Pre-trained Transformer}，``生成式预训练 Transformer''，名字里最后一个词就是它。想弄明白现在的 LLM 在做什么，Transformer 是绕不开的第一站。市面上的讲解不少，但大多要么照着论文堆公式，要么停在最浅的比喻层面，于是我想自己写一篇，把它的结构和计算流程从头到尾走一遍，也方便日后忘记时回头翻翻。

为了让计算能摊在纸上手算，例子里只用四个 token、四维表示这样的尺寸；完整问答示例换成八个 token、十二维表示。两者都是教学尺寸，不是原论文的模型配置；优化器调度、分词算法和推理加速等工程话题也不展开，先把最核心的计算链条接通。本文讨论的是 \textit{Attention Is All You Need} 中的原始编码器—解码器 Transformer，采用图中``先 Add、后 Norm''的连接方式，侧重直观理解，省略了部分纯计算性的细节，并附上了笔者自己的理解，如有错误，还请指正；主要参考了原论文、《动手学深度学习》的相关章节与 PyTorch 文档。

\section{序列任务与注意力}

\subsection{序列任务的困难}

Transformer 最初是为机器翻译设计的，所以先从这类任务的难处说起。机器翻译、问答等任务都是从一个序列产生另一个序列：比如输入``谁是第一个登上月球的人？''，输出一个人名。这里光认识单个词还不够——``第一个''``月球''``人''要组合在一起，才约束住问题的含义。也就是说，模型必须让不同位置的信息发生交互。

一种自然的处理方式是顺序读取。循环神经网络（RNN）在第 $i$ 个位置计算
\begin{equation}
\bm h_i=f_\theta(\bm h_{i-1},\bm x_i).
\label{eq:rnn}
\end{equation}
这里 $\bm h_i$ 是读到当前位置时的内部状态。由于计算 $\bm h_i$ 需要先得到 $\bm h_{i-1}$，同一个序列内部存在串行依赖；相隔很远的位置之间，也往往需要经过多步状态传递。LSTM 等结构改善了记忆问题，但没有消除这种逐位置计算的依赖。

注意力提供了另一种办法：处理当前位置时，直接对允许访问的位置打分，再汇总其中的信息。它在 Transformer 之前就已经被用于神经机器翻译；Transformer 的关键并不是首次发明注意力，而是用注意力承担主要的序列信息交互，不再依赖原来的循环结构。

\begin{figure}[htbp]
\centering
\begin{tikzpicture}[x=1cm,y=1cm]
\node[flowlab,anchor=east] at (-1.15,0.9) {RNN};
\foreach \i/\x in {1/0,2/3,3/6,4/9}{
  \node[flowbox,minimum width=1.8cm,minimum height=0.75cm] (r\i) at (\x,0.9) {$\bm h_{\i}$};
}
\draw[flowarr] (r1)--(r2); \draw[flowarr] (r2)--(r3); \draw[flowarr] (r3)--(r4);
\node[flowlab,anchor=east] at (-1.15,-1.0) {Attention};
\foreach \i/\x in {1/0,2/3,3/6,4/9}{
  \node[flowbox,minimum width=1.8cm,minimum height=0.75cm] (a\i) at (\x,-1.0) {$\bm x_{\i}$};
}
\node[flowbox,minimum width=3.0cm,minimum height=0.75cm] (z) at (4.5,-2.5) {第 4 个位置的新特征};
\foreach \i in {1,2,3,4}{\draw[flowarr] (a\i.south)--(z.north);}
\end{tikzpicture}
\caption{RNN 的串行依赖与注意力的直接汇总}
\label{fig:motivation}
\end{figure}

需要说明的是，注意力并没有``根治''长距离依赖，Transformer 也不是在任何情况下都更快：完整注意力要考虑所有位置的配对，序列很长时计算量依然可观；层与层之间仍要顺序执行，生成答案时也还是一个 token 一个 token 地往外产生。它真正带来的，是一条更直接的信息通路，以及同一层内天然适合并行的计算方式。

\subsection{输入与输出}

原始 Transformer 采用 sequence-to-sequence（seq2seq）形式：Encoder 把输入序列转换成一串连续向量；Decoder 读取这些向量以及已知的输出前缀，预测下一个 token。输入序列与输出序列不必一样长。

``seq2seq''描述输入、输出之间的组织方式，``注意力''描述内部如何汇总信息，两者并不冲突。所有向量默认按行摆放；$n$ 表示输入长度，$t$ 表示当前输出前缀长度，$d$ 表示每个位置的表示维度。暂时不写 batch 维度，以免它遮住核心计算。

\subsection{整体结构}

在拆开每个部件之前，先看一眼全景（图~\ref{fig:architecture}）。原始 Transformer 由两座``塔''组成：左半 Encoder 从下往上是输入嵌入、位置编码和若干个相同的编码层；右半 Decoder 结构类似，但每个层里多出一个接收 Encoder 输出的子模块。现在看不懂细节没关系——底部圆圈与``$+$''表示加入位置编码，绕过各子模块的连线是残差相加，这些后面都会逐个拆开。把这张图记在心里，接下来沿着数据流动的顺序，从进入模型的第一步说起：一串文字是怎么变成向量的。

\begin{figure}[htbp]
\centering
\includegraphics[width=0.40\textwidth]{{Overview of transformer}.png}
\caption{原始 Transformer 结构（原论文 Figure 1）}
\label{fig:architecture}
\end{figure}

\section{输入表示}

\subsection{Token 与词嵌入}

计算机没法直接处理文字，第一步是把文字拆成基本单位并编号。Token 就是这样的基本单位——由分词器切分得到，未必是完整单词，也可能是字或子词。为方便手算，这里把 \texttt{I love deep learning} 当作四个 token。词表为它们分配编号；编号只负责索引，不表示语义大小。

Embedding（嵌入）是一张可训练的查找表。设词表为 $\mathcal V$，则嵌入表
\begin{equation}
\bm B\in\R^{|\mathcal V|\times d}
\label{eq:embeddingtable}
\end{equation}
为每个 token 存放一个 $d$ 维向量。输入一个编号，就取出对应的一行；这不是把编号本身当成数值特征参与语义计算。

例如，假设查表结果如表~\ref{tab:embedding} 所示。所有数字都是人为选择的演示值，不代表真实模型中的词义。
\begin{table}[htbp]
\centering
\caption{四个 token 的示意嵌入}
\label{tab:embedding}
\renewcommand{\arraystretch}{1.2}
\begin{tabular}{ccc}
\toprule
Token & 编号 & 查表得到的向量\\
\midrule
\texttt{I} & 12 & $[1,0,1,0]$\\
\texttt{love} & 37 & $[0,1,1,0]$\\
\texttt{deep} & 81 & $[0,2,0,1]$\\
\texttt{learning} & 96 & $[1,1,0,1]$\\
\bottomrule
\end{tabular}
\end{table}

按句子顺序把取出的四行叠起来，就得到当前句子的嵌入矩阵
\begin{equation}
\bm E_s=
\begin{bmatrix}
1&0&1&0\\
0&1&1&0\\
0&2&0&1\\
1&1&0&1
\end{bmatrix}\in\R^{4\times4}.
\label{eq:embeddingexample}
\end{equation}
注意区分：$\bm B$ 是整张词表的参数；$\bm E_s$ 是当前这句话查表得到的中间结果。一个位置对应一行，一个特征对应一列。

这些向量并不是人工定义好的``词义坐标''。从头训练时可以随机初始化，随后通过最终任务的损失与其他参数一起调整。每个维度通常没有预先规定的人类语义。参数固定时，同一个 token 的初始查表结果相同；经过后面的上下文计算后，它在不同句子中的表示才可以不同。

\subsection{位置编码}

有了词的向量，还缺顺序信息。``人咬狗''和``狗咬人''用的是同一组词，含义却完全相反。而后面的自注意力按内容汇总信息，本身并不感知``谁在前谁在后''，所以顺序必须显式地加进向量里。

原始 Transformer 生成一个与嵌入矩阵同形状的位置矩阵 $\bm P$，再相加。严格写出原论文中的嵌入缩放，有
\begin{equation}
\bm X^{(0)}=\sqrt d\,\bm E_s+\bm P,
\qquad \bm E_s,\bm P,\bm X^{(0)}\in\R^{n\times d}.
\label{eq:positionadd}
\end{equation}
例如，八个 token、每个十二维时，相加前后都是 $8\times12$，不是 $8\times13$。两个矩阵是在相同位置上逐元素相加，不是首尾拼接。

原始的正弦位置编码为
\begin{align}
P_{p,2j}&=\sin\!\left(\frac{p}{10000^{2j/d}}\right),\notag\\
P_{p,2j+1}&=\cos\!\left(\frac{p}{10000^{2j/d}}\right),
\label{eq:sinusoidal}
\end{align}
其中位置 $p$ 从零开始；对偶数 $d$，有 $j=0,\ldots,d/2-1$。不同维度使用不同频率，因此一个位置对应一整组有规律的数，而不是单独附加一个整数编号。至于为什么偏偏用正弦和余弦：原文的一个动机是这组波形有个好性质——对任意固定偏移 $k$，位置 $p+k$ 的编码都能写成位置 $p$ 编码的线性函数，方便模型学到相对位置关系。这里不必深究，把它理解成``每个位置对应一串固定规律的数''就够了。

\subsection{相加不会混淆吗？}

到这里，一个疑问很自然地冒出来：内容和位置是直接加在一起的，会不会有两个不同的``token + 位置''组合，加出同一个向量，让模型无从分辨？

如果这些数字是任意选的，确实可能。把式~\eqref{eq:positionadd} 的共同缩放暂时吸收入 $\bm e$ 的记号，取
\begin{align*}
\bm e(a)&=[0.1,0.2,0],&\bm e(b)&=[0,0.1,-0.1],\\
\bm p_1&=[0.1,0.1,0.1],&\bm p_2&=[0,0,0],
\end{align*}
就有
\begin{equation}
\bm e(a)+\bm p_2=\bm e(b)+\bm p_1=[0.1,0.2,0].
\label{eq:collision}
\end{equation}
单看加出来的结果，区分不出它来自哪个组合。

那 Transformer 为什么还敢直接相加？因为相加的双方并不是随机的数字：位置编码是按固定规律生成的波形（式~\eqref{eq:sinusoidal}），嵌入是从训练中学出来的参数。模型从训练一开始就在``内容 + 位置''的联合表示上工作，完全可以学出适合任务的组合方式，并不需要先单独还原出``纯净的词义''、再单独还原出位置。式~\eqref{eq:collision} 说明的只是这种写法没有数学上的唯一性保证，而不是说相加一定会丢失信息。

到这里，输入侧的准备就完成了：每个 token 都带着内容与位置信息，变成了一行向量。接下来进入 Transformer 最核心的部分——注意力。

\section{自注意力：从单头到多头}
\label{sec:singleattention}

\subsection{注意力层的输出}

先澄清一个容易产生的直觉偏差：注意力并不是从句子里挑出一个``最重要的词''。它为每个查询位置分别生成一份上下文汇总——第 $i$ 个位置有自己的分配比例，决定从各个位置分别取多少信息，而且这份比例随输入内容改变。

为了能把整个计算搬到纸上手算，下面直接给注意力子层指定一个输入
\begin{equation}
\bm X=
\begin{bmatrix}
1&0&1&0\\
0&1&1&0\\
0&2&0&1\\
1&1&0&1
\end{bmatrix}.
\label{eq:toyX}
\end{equation}
四行仍对应 \texttt{I/love/deep/learning}。它的数值恰好沿用式~\eqref{eq:embeddingexample} 的嵌入矩阵，但请不要把它当成``嵌入加上位置编码''之后的结果——手算部分不再叠加位置编码，也省略投影偏置，专门用一个干净的输入把流程走通。

\subsection{Query、Key、Value 的来源}

同一个位置要参与两类工作：一类是决定``如何匹配和分配''，另一类是提供``真正被汇总的内容''。Query 与 Key 参与匹配，Value 参与最后的汇总。三者由输入经过不同的可训练线性投影得到：
\begin{equation}
\bm Q=\bm X\bm W_Q,\qquad
\bm K=\bm X\bm W_K,\qquad
\bm V=\bm X\bm W_V.
\label{eq:qkv}
\end{equation}
这里 $\bm X$ 不是 $\bm Q$，而是产生 $\bm Q,\bm K,\bm V$ 的原材料。

假设选择下面三套 $4\times2$ 的权重：
\begin{equation}
\bm W_Q=\begin{bmatrix}1&0\\0&1\\0&1\\0&-1\end{bmatrix},\quad
\bm W_K=\begin{bmatrix}0&0\\1&0\\0&1\\0&0\end{bmatrix},\quad
\bm W_V=\begin{bmatrix}1&0\\0&1\\1&0\\0&1\end{bmatrix}.
\label{eq:toyweights}
\end{equation}
``乘权重矩阵''就是按系数组合原来的特征。对于 $\bm x=[a,b,c,d]$，这三个矩阵分别实现
\begin{equation}
\bm q=[a,b+c-d],\qquad
\bm k=[b,c],\qquad
\bm v=[a+c,b+d].
\label{eq:projectionmeaning}
\end{equation}
真实模型不会固定使用这几条规则，而是在训练中调整系数。对所有行应用同一套规则，得到
\begin{equation}
\bm Q=\begin{bmatrix}1&1\\0&2\\0&1\\1&0\end{bmatrix},\quad
\bm K=\begin{bmatrix}0&1\\1&1\\2&0\\1&0\end{bmatrix},\quad
\bm V=\begin{bmatrix}2&0\\1&1\\0&3\\1&2\end{bmatrix}.
\label{eq:toyqkv}
\end{equation}
三者形状都是 $4\times2$。到这里还没有交换 token 之间的信息，只是每一行分别做了特征变换。

\subsection{分数矩阵}

先只看第四个位置 \texttt{learning}。它的 Query 是 $\bm q_4=[1,0]$，与各 Key 点积：
\begin{align}
\bm q_4\bm k_1^\top&=0,&
\bm q_4\bm k_2^\top&=1,\notag\\
\bm q_4\bm k_3^\top&=2,&
\bm q_4\bm k_4^\top&=1.
\label{eq:fourthscores}
\end{align}
得到 $[0,1,2,1]$，即 \texttt{learning} 对四个来源的匹配分数。注意分数里包含它自己——自注意力允许位置查看自己，但自己的分数并不保证最高：这里 \texttt{learning} 与 \texttt{deep} 的匹配分数（$2$）就高于与自己（$1$）。

其他查询位置也进行相同计算。将它们集中成矩阵乘法，就是
\begin{equation}
\bm S=\bm Q\bm K^\top
=\begin{bmatrix}
1&2&2&1\\
2&2&0&0\\
1&1&0&0\\
0&1&2&1
\end{bmatrix}.
\label{eq:scorematrix}
\end{equation}
转置的作用是把每个 Key 从一行变成一列，使矩阵乘法实现所有 Query 与所有 Key 的点积：
\[
(4\times2)(2\times4)\longrightarrow4\times4.
\]
因此 $S_{ij}$ 的行 $i$ 是接收信息的位置，列 $j$ 是候选来源。$\bm S$ 的列已经不是原始特征维度，而是序列中的位置。

\subsection{缩放与 softmax}

这个头的 Query 与 Key 都是二维，所以 $d_k=2$。先把每个分数除以 $\sqrt2$；矩阵形状不变。第四行变成
\[
[0,1,2,1]/\sqrt2\approx[0,0.707,1.414,0.707].
\]
为什么要缩放？在一个简化假设下，Query 和 Key 的各分量独立、均值为零、方差为一，则
\begin{equation}
\operatorname{Var}\!\left(\sum_{r=1}^{d_k}q_rk_r\right)=d_k.
\label{eq:variance}
\end{equation}
点积的标准差随 $\sqrt{d_k}$ 增长——维度一高，分数动辄拉开很大差距，softmax 就会过度集中在少数位置上。提前除以 $\sqrt{d_k}$，相当于把这个尺度拉回来。当然，训练后的分量未必真满足上面的独立假设，$\sqrt{d_k}$ 更像一个稳妥的工程选择。

接下来对每一行分别做 softmax：
\begin{equation}
A_{ij}=\frac{\exp(S_{ij}/\sqrt{d_k})}
{\sum_{r=1}^{n}\exp(S_{ir}/\sqrt{d_k})}.
\label{eq:rowsoftmax}
\end{equation}
第四行先取指数约为 $[1,2.028,4.113,2.028]$，再除以这一行的总和，得到
\begin{equation}
\bm a_4\approx[0.109,0.221,0.449,0.221].
\label{eq:fourthweights}
\end{equation}
四个数相加为一，表示当前查询位置如何分配汇总比例。所有行一起算得
\begin{equation}
\bm A=\softmax(\bm S/\sqrt2)
\approx\begin{bmatrix}
0.165&0.335&0.335&0.165\\
0.402&0.402&0.098&0.098\\
0.335&0.335&0.165&0.165\\
0.109&0.221&0.449&0.221
\end{bmatrix}.
\label{eq:attentionmatrix}
\end{equation}
$\bm A$ 叫作注意力权重矩阵，每一行是一种分配比例。它随输入内容和参数变化——同一句话换掉一个词，权重可能就是另一番样子。矩阵也不必对称：第 $i$ 行第 $j$ 列说的是``位置 $i$ 从 $j$ 取多少''，第 $j$ 行第 $i$ 列说的是反向的取法，本来就是两件事。

\subsection{加权求和输出}

第四个位置的输出是四个 Value 的加权和：
\begin{align}
\bm z_4&=\sum_{j=1}^{4}A_{4j}\bm v_j\notag\\
&\approx0.109[2,0]+0.221[1,1]
       +0.449[0,3]+0.221[1,2]\notag\\
&\approx[0.660,2.009].
\label{eq:weightedvalues}
\end{align}
以第一维为例：$0.109\times2+0.221\times1+0.449\times0+0.221\times1\approx0.66$。式中的权重只保留了三位小数，最终结果按未舍入的权重计算。

这和平均数很像，但不是每个来源都占 $1/4$，而是每个查询位置根据内容决定自己的权重。对所有位置一次完成汇总，就是
\begin{equation}
\bm Z=\bm A\bm V
\approx\begin{bmatrix}
0.830&1.670\\
1.304&0.891\\
1.170&1.160\\
0.660&2.009
\end{bmatrix}\in\R^{4\times2}.
\label{eq:attentionout}
\end{equation}
$\bm A$ 是``怎么取''的规则，$\bm Z$ 是``实际取到了什么''的特征。前者每行四个权重，后者每行两个输出特征，不能混为一谈。

现在把三步连起来，得到熟悉的公式：
\begin{equation}
\boxed{\operatorname{Attention}(\bm Q,\bm K,\bm V)
=\softmax\!\left(\frac{\bm Q\bm K^\top}{\sqrt{d_k}}\right)\bm V.}
\label{eq:attention}
\end{equation}
\begin{figure}[htbp]
\centering
\includegraphics[width=0.40\textwidth]{{Scaled Dot-Product Attention}.png}
\caption{缩放点积注意力（原论文 Figure 2 左）}
\label{fig:scaledattention}
\end{figure}

整个公式其实就是三步：$\bm Q\bm K^\top$ 完成匹配，softmax 把分数变成分配比例，再按比例汇总 $\bm V$ 的内容。后面要讲的多头注意力、交叉注意力，变的都是这三步里各矩阵的来源，套路本身不再变。

\subsection{多头注意力}

恢复到 $\bm X\in\R^{8\times12}$ 的教学尺寸。一个头只有一套注意力分配方式：对某个来源给出权重 $a$，就会把该来源的整个 Value 向量乘上同一个 $a$。多头允许同时计算多套分配与内容表示，再合并结果。

设有三个头，每个头的 Query、Key、Value 都是四维。第 $h$ 个头计算
\begin{equation}
\bm Q_h=\bm X\bm W_Q^{(h)},\quad
\bm K_h=\bm X\bm W_K^{(h)},\quad
\bm V_h=\bm X\bm W_V^{(h)},
\label{eq:headprojection}
\end{equation}
其中每个权重矩阵都是 $12\times4$。因此每个头都能组合完整的十二维输入，并不是先截取原向量的前四维、中间四维和后四维。它们各自得到 $8\times8$ 的权重矩阵，以及 $8\times4$ 的输出 $\bm Z_h$。

需要说明的是，并不存在``一个头专门负责主谓搭配、一个头专门负责指代''这样的预设分工。各头只是参数不同，靠同一个损失联合训练；它们最终学出怎样的分工，是训练的结果，不是设计时的规定。

\begin{figure}[htbp]
\centering
\includegraphics[width=0.50\textwidth]{{Multi-Head Attention}.png}
\caption{多头注意力（原论文 Figure 2 右）}
\label{fig:multihead}
\end{figure}

一般写成
\begin{equation}
\MHA(\bm X)=\Concat(\bm Z_1,\ldots,\bm Z_H)\bm W_O,
\quad \bm W_O\in\R^{Hd_v\times d}.
\label{eq:multihead}
\end{equation}
实现中可以把多个投影合成一次大矩阵乘法，再把投影后的结果拆成头。例如
\[
\bm X[\bm W_Q^{(1)}\mid\bm W_Q^{(2)}\mid\bm W_Q^{(3)}]
=[\bm Q_1\mid\bm Q_2\mid\bm Q_3].
\]
所谓``拆头''常常发生在这一步，而不是在原始输入上直接切掉部分信息。

\subsection{输出维度}

各头的内部宽度并不必须是四。完全可以让每个头输出 $8\times5$，三个头拼成 $8\times15$，再选择
\begin{equation}
\bm W_O\in\R^{15\times12},\qquad
(8\times15)(15\times12)\longrightarrow8\times12.
\label{eq:alternativeheads}
\end{equation}
回到十二维，是为了与网络主路径的宽度一致，方便残差相加和逐层连接。常见的 $d_v=d/H$ 只是一种让内部总宽度、计算量和参数量都比较可控的取法，并非硬性规定。

八行之所以保留，是因为这里有八个 Query，需要分别输出八个向量。行数由查询位置数决定，列数则由特征设计和输出投影决定。

至此，注意力的两种形态——单头与多头——都齐了。不过光有注意力还组不成一个完整的层：把它装进网络，还需要残差连接、归一化和前馈网络。

\section{Add、Norm 与 FFN}
\label{sec:block}

\subsection{残差连接}

暂时不考虑归一化，用 $F_\theta$ 表示注意力或 FFN 这样的可训练模块。没有残差时，模块输出直接作为结果；有残差时，再加上该模块的输入：
\begin{equation}
\text{无残差：}\quad \hat{\bm y}=F_\theta(\bm x),
\qquad
\text{有残差：}\quad \hat{\bm y}=\bm x+F_\theta(\bm x).
\label{eq:residual}
\end{equation}
残差把输入分成两路：一路经过模块变换，一路原样直通，最后相加。注意，直通的那一路从训练的第一步就在起作用——损失始终在评判 $\bm x+F_\theta(\bm x)$ 这整个整体，而不是先把 $F_\theta$ 单独训好、最后再把 $\bm x$ 补上去。

\begin{figure}[htbp]
\centering
\begin{tikzpicture}[x=1cm,y=1cm]
\node[flowlab] (in) at (0,0) {$\bm x$};
\node[flowbox,minimum width=2.8cm,minimum height=1cm] (f) at (3.0,0) {$F_\theta(\bm x)$\\注意力或 FFN};
\node[sumcircle] (add) at (5.8,0) {$+$};
\node[addnorm,minimum width=1.8cm,minimum height=1cm] (ln) at (8.1,0) {LayerNorm};
\node[flowlab] (out) at (10.4,0) {$\bm y$};
\draw[flowarr] (in)--(f); \draw[flowarr] (f)--(add);
\draw[flowarr] (add)--(ln); \draw[flowarr] (ln)--(out);
\draw[flowarr] (0.9,0)--(0.9,1.35)--node[above,font=\small]{直接路径：同一个子模块的输入 $\bm x$}(5.8,1.35)--(add.north);
\fill[gray!75] (0.9,0) circle(1.3pt);
\end{tikzpicture}
\caption{残差连接与归一化}
\label{fig:residual}
\end{figure}

例如，我们希望实现 $T(x)=3x$，分支取 $F_\theta(x)=\theta x$。没有残差时需要 $\theta=3$；有残差时只需要 $\theta=2$，因为 $x+2x=3x$。目标没有改变，达到目标所需的参数改变了。

取训练样本 $x=2$、目标 $y^*=6$。有残差时，真正参与损失的是
\begin{equation}
L(\theta)=(2+2\theta-6)^2,
\qquad \frac{\mathrm dL}{\mathrm d\theta}=4(2+2\theta-6).
\label{eq:residualloss}
\end{equation}
最优解是 $\theta=2$：分支输出 $2\times2=4$，加上输入的 $2$ 恰好等于 $6$。如果让分支直接去拟合目标 $6$，加回输入反而得到 $8$。也就是说，残差结构下分支要学的不是目标本身，而是``距离目标还差多少''。

用 $T(\bm x)$ 表示理想变换，那么 $F_\theta(\bm x)$ 学的就是差值 $T(\bm x)-\bm x$，也就是``更新量''。不过要留意：训练数据只给最终答案，中间并没有一份现成的 $T(\bm x)$ 摆在那里供相减，网络是从整体损失中自己摸索出每层该输出什么更新的。

\keypoint{残差分支学的是``更新量''，但更新量并不总是很小。}

从 $2$ 变到 $6$，需要的更新就是 $4$，不小也不该小；只有某个子层确实不需要改变输入时，零更新才合理。当然，``能学差值''只是代数上的可能性，能不能学好仍取决于网络的表达能力和训练过程。

\subsection{直接路径与反向传播}

残差对训练的好处，从反向传播的角度看得更清楚。先用一维说明：设 $y=x+F_\theta(x)$，后面的损失为 $L$，链式法则给出
\begin{equation}
\frac{\partial L}{\partial x}
=\frac{\partial L}{\partial y}\bigl(1+F_\theta'(x)\bigr).
\label{eq:residualgradient}
\end{equation}
没有残差时，括号里只有 $F_\theta'(x)$。假如这个导数是 $0.01$、后一层传来的梯度是 $1$，两种情况分别得到 $0.01$ 和 $1.01$——哪怕分支学得不好、导数很小，梯度仍能沿着那个 $+1$ 的直通路回到浅层。这不是``永不消失''的保证，但确实给了梯度一条不依赖分支质量的退路。

多维情况对应 $\bm I+\bm J_F$，其中 $\bm J_F$ 是分支的雅可比矩阵。原始 Transformer 后面还接 LayerNorm，梯度也要经过它，不能把整层视为完全不变的恒等通道。只把分支设成零，原图的子层仍输出 $\LN(\bm x)$，而不是严格的 $\bm x$。

\subsection{LayerNorm 的归一化对象}

先做
\begin{equation}
\bm U=\MHA(\bm X),\qquad \bm Z=\bm X+\bm U,
\qquad \bm Y=\LN(\bm Z).
\label{eq:lnexpanded}
\end{equation}
把前两项代入，就是常见的写法 $\bm Y=\LN(\bm X+\MHA(\bm X))$。写这么长只是为了看清结构——LN 实际接收的只有相加后的 $\bm Z$ 这一个矩阵。

对某一行 $\bm z=[z_1,\ldots,z_d]$，计算
\begin{equation}
\mu=\frac1d\sum_{j=1}^{d}z_j,\qquad
\sigma^2=\frac1d\sum_{j=1}^{d}(z_j-\mu)^2,
\label{eq:lnstats}
\end{equation}
再逐维得到
\begin{equation}
\LN(\bm z)_j=\gamma_j\frac{z_j-\mu}{\sqrt{\sigma^2+\epsilon}}+\beta_j.
\label{eq:layernorm}
\end{equation}
$\epsilon$ 是避免数值问题的小常数；$\bm\gamma,\bm\beta$ 是可训练的缩放与平移参数。每个 token 分别沿自己的特征维度计算统计量；它不把整句话混在一起平均，也不是沿 batch 计算。

例如 $\bm z=[1,2,3,4]$ 的均值为 $2.5$、方差为 $1.25$。暂取 $\bm\gamma=\bm1,\bm\beta=\bm0$ 并忽略很小的 $\epsilon$，得到
\[
\LN(\bm z)\approx[-1.342,-0.447,0.447,1.342].
\]
若把整个向量放大一百倍，换成 $[100,200,300,400]$，标准化后的结果几乎不变——LN 关心的是各维度之间的相对起伏，而不是绝对尺度。当然它也只是调整尺度，并不创造信息；配合可训练的 $\bm\gamma$ 与 $\bm\beta$，最终输出的均值和方差也不必再是零和一。

尺度为什么值得管？忽略偏置，若输入整体放大 $c$ 倍，投影的 $\bm Q,\bm K$ 也随之放大，点积变成 $c^2\bm Q\bm K^\top$，softmax 的集中程度就可能明显改变。LN 有助于让后续层面对较受控的数值尺度，并改善深层网络的训练性质。

原图使用先相加、再归一化的 Post-LN。把 LN 放到分支前的 Pre-LN 是另一种结构，例如 $\bm y=\bm x+F_\theta(\LN(\bm x))$；两者不是等价的括号挪动，训练性质也不同。本文只追踪原图，不在这里展开变体。

\subsection{FFN 与非线性}

注意力负责按内容汇总不同位置的信息；FFN 则对每个位置已经得到的特征继续做变换。原始结构为
\begin{equation}
\FFN(\bm x)=\ReLU(\bm x\bm W_1+\bm b_1)\bm W_2+\bm b_2.
\label{eq:ffn}
\end{equation}
例如可选 $12\rightarrow48\rightarrow12$ 的宽度，其中 $\ReLU(x)=\max(x,0)$ 是逐元素取正部。中间扩展提供更多候选特征，非线性使它不再只是一次固定的线性组合。同一层的各位置共用 FFN 参数，但分别计算。

一个极小的例子就能展示用意。让输入为一个数 $x$，两次线性变换和一次 ReLU 实现
\begin{equation}
 x\longrightarrow[x,-x]
 \longrightarrow[\max(x,0),\max(-x,0)]
 \longrightarrow |x|.
\label{eq:ffnabs}
\end{equation}
输入 $-2$，依次得到 $[-2,2]$、$[0,2]$，最后得到 $2$。这正是一次 $1\rightarrow2\rightarrow1$ 的 FFN。

如果拿掉非线性、暂不写偏置，则
\[
(\bm x\bm W_1)\bm W_2=\bm x(\bm W_1\bm W_2),
\]
两次线性变换合并成一次，包含偏置时也只是合并成一个仿射变换——中间维度升得再高，表达能力也不会变。升维的价值，要靠非线性来兑现。至于注意力里也有 softmax 这样的非线性，两者并不冲突：FFN 的角色是提供逐位置的非线性特征变换，在这一步不重新混合 token 之间的信息。

\subsection{完整的 Encoder Block}

忽略 dropout，原始 Encoder 一层只需写成
\begin{align}
\bm H&=\LN_1\bigl(\bm X+\MHA(\bm X)\bigr),\notag\\
\bm Y&=\LN_2\bigl(\bm H+\FFN(\bm H)\bigr).
\label{eq:encoderblock}
\end{align}
每个主要子模块都配一组 Add 与 Norm，所以一层有两组。第二次 Add 加回的是 FFN 的输入 $\bm H$，不是最初的 $\bm X$。两个 LN 有各自的可训练参数，但输出形状都不变。

原论文训练时还在分支输出等位置使用 dropout，更完整的第一行是 $\LN_1(\bm X+\operatorname{Dropout}(\MHA(\bm X)))$。本文为突出主线，后文一律省略它。

到这里，Encoder 的一层就完整了。把 $N$ 块这样的结构堆起来，就得到整个 Encoder：
\begin{equation}
\bm M=\operatorname{Encoder}(\bm X^{(0)})\in\R^{n\times d}.
\label{eq:memory}
\end{equation}
各层结构相同，但参数不同——后一层重新计算自己的 Query、Key、Value 和注意力，不是重复使用第一层的权重矩阵。$\bm M$ 的每一行仍对应问题中的一个位置，只是已经受到上下文与多层特征变换的影响。

值得强调的是，Encoder 的输出是一整串向量，而不是把整句话压缩成的一个向量——更早的 seq2seq 系统正是后一种做法，句子一长，一个向量就装不下全部信息。Transformer 把这串表示原封不动地交给 Decoder，让它按需取用。下一节就看 Decoder 怎么接住它。

\section{Decoder：掩码自注意力与交叉注意力}
\label{sec:decoder}

Encoder 侧的零件到这里都齐了。Decoder 的每一层用的也是这些零件，只有两处不同：自注意力带上了因果掩码，中间多出一个读取 $\bm M$ 的交叉注意力。回到图~\ref{fig:architecture} 的右半边，从下往上看。

\subsection{掩码自注意力}

第一次生成时，Decoder 的输入只有开始标记 \bos。以后则是 \bos\ 加上已经生成的 token。图中右下角的 Output Embedding，表示把这串已知输出转成向量，不是把向量变成最终文字。

令当前前缀的表示为 $\bm U\in\R^{t\times d}$。在掩码自注意力中，$Q,K,V$ 都来自 $\bm U$，但第 $i$ 个位置只能访问 $j\le i$ 的位置。用一个加在分数上的矩阵表示规则：
\begin{equation}
C_{ij}=\begin{cases}
0,&j\le i,\\
-\infty,&j>i.
\end{cases}
\qquad
\bm A=\softmax\!\left(\frac{\bm Q\bm K^\top}{\sqrt{d_k}}+\bm C\right).
\label{eq:causalmask}
\end{equation}
对于三个前缀 token，$\bm C$ 为
\[
\begin{bmatrix}
0&-\infty&-\infty\\
0&0&-\infty\\
0&0&0
\end{bmatrix}.
\]
因为 $\exp(-\infty)=0$，不允许访问的未来位置得到零权重。掩码不是把当前输入 token 自己隐藏起来，而是阻止它提前读取后面的 token。

这一模块后面仍有 Add 与 Norm。掩码自注意力的结果已经融入答案前缀内部的信息，但要根据问题作答，还需要下一步。

\subsection{交叉注意力}

设右边经过前一子模块后得到 $\bm H$。交叉注意力对它和 Encoder 输出分别投影：
\begin{equation}
\bm Q=\bm H\bm W_Q^{\mathrm{cross}},\qquad
\bm K=\bm M\bm W_K^{\mathrm{cross}},\qquad
\bm V=\bm M\bm W_V^{\mathrm{cross}}.
\label{eq:crossqkv}
\end{equation}
Query 来自当前的答案表示；Key、Value 来自问题表示。公式仍然是式~\eqref{eq:attention}，改变的是输入来源。

例如，问题有八个 token，答案前缀有三个 token，每个头使用四维 Query、Key、Value。那么
\begin{equation}
\underbrace{\bm Q}_{3\times4}
\underbrace{\bm K^\top}_{4\times8}
\longrightarrow
\underbrace{\bm A}_{3\times8},
\qquad
\underbrace{\bm A}_{3\times8}
\underbrace{\bm V}_{8\times4}
\longrightarrow
\underbrace{\bm Z}_{3\times4}.
\label{eq:crossshape}
\end{equation}
这表示答案的三个位置各自从问题的八个位置汇总信息。答案短、问题长并不冲突：Query 只决定每个答案位置``去问题里各取多少''，两边的长度本来就是相互独立的。

\begin{table}[htbp]
\centering
\caption{三种注意力的连接关系（尺寸按单个头计）}
\label{tab:threeattentions}
\small\renewcommand{\arraystretch}{1.35}
\begin{tabularx}{\textwidth}{@{}>{\raggedright\arraybackslash}p{3.3cm}YYc@{}}
\toprule
模块 & Query 的来源 & Key、Value 的来源 & 权重矩阵\\
\midrule
Encoder 自注意力 & 问题的当前表示 & 同一份问题表示 & $8\times8$\\
Decoder 掩码自注意力 & 答案前缀的当前表示 & 同一份前缀表示，屏蔽未来 & $3\times3$\\
Decoder 交叉注意力 & 右侧前一子模块的输出 & Encoder 的最终输出 & $3\times8$\\
\bottomrule
\end{tabularx}
\end{table}

\subsection{完整的 Decoder Block}

忽略 dropout，一个 Decoder Block 完整写成
\begin{align}
\bm H&=\LN_1\bigl(\bm U+\operatorname{MaskedMHA}(\bm U)\bigr),\notag\\
\bm R&=\LN_2\bigl(\bm H+\operatorname{CrossMHA}(\bm H,\bm M)\bigr),\notag\\
\bm D&=\LN_3\bigl(\bm R+\FFN(\bm R)\bigr).
\label{eq:decoderblock}
\end{align}
三个主要子模块对应三组 Add 与 Norm，每一步都保持右侧 $t\times d$ 的形状。第二行加回的是右侧的 $\bm H$，不是左侧的 $\bm M$。

把 Decoder Block 和 Encoder Block 放在一起对比：自注意力换成了带因果掩码的版本，中间多出一个读取 $\bm M$ 的交叉注意力，其余的残差、LN、FFN 完全一致。下一层 Decoder 继续处理 $\bm D$，仍可访问同一份 $\bm M$；各层使用自己的投影参数。

\section{逐步生成与训练}
\label{sec:generation}

\subsection{一个完整例子}

结构讲完了，让整个模型完整跑一遍。假设任务是问答：输入``谁是第一个登上月球的人？''，期望模型简短地回答``阿姆斯特朗''（第一位踏上月球的 Neil A. Armstrong）。当然，能答对的前提是接受过相关训练——只搭好结构、还没训练的网络，不会仅凭框架图就知道这条事实。

人为把问题分为八个 token：
\[
[\text{谁， 是， 第一个， 登上， 月球， 的， 人， ？}].
\]
这只是教学切分。Encoder 产生 $\bm M\in\R^{8\times12}$；为了展示方便，答案按单字逐个生成。

假设当前前缀是``\bos\ 阿姆''，它有三个 token。右边完成式~\eqref{eq:decoderblock} 并经过全部 Decoder 层，得到一个 $3\times12$ 的矩阵。每行都是对应位置的上下文化表示。生成下一个 token 时，取最后一行，也就是当前``姆''所在位置的向量：
\begin{equation}
\bm d_{\mathrm{last}}\in\R^{1\times12}.
\label{eq:lastvector}
\end{equation}
取最后一行并不表示忽略前面位置；这个向量已经通过注意力读取了允许访问的前缀和问题信息。

\subsection{Linear 与 Softmax}

假设目标词表有 $10{,}000$ 个 token。输出线性层采用
\begin{equation}
\bm W_{\mathrm{vocab}}\in\R^{12\times10000},\qquad
\bm\ell=\bm d_{\mathrm{last}}\bm W_{\mathrm{vocab}}+\bm b
\in\R^{1\times10000}.
\label{eq:vocablinear}
\end{equation}
这里每个分数对应词表中的一个候选 token，而不是一个输入位置。再计算
\begin{equation}
P(y_t=j\mid s,y_{<t})=
\frac{\exp(\ell_j)}{\sum_{r=1}^{|\mathcal V_{\mathrm{out}}|}\exp(\ell_r)}.
\label{eq:vocabsoftmax}
\end{equation}
于是得到了下一个 token 的概率分布。

这和注意力里的 softmax 是同一种数学运算，但对象不同：注意力在``信息来源的位置''之间分配权重，输出层在``词表候选''之间分配概率。softmax 本身不负责选中某个 token；还需要解码规则，例如取概率最高者，或按照分布采样。

为展示最简单的流程，本文采用逐步取最大概率的贪心选择：
\[
\hat y_t=\argmax_j P(y_t=j\mid s,\hat y_{<t}).
\]
这只是最简单的演示选择：训练后的模型未必给正确答案最高概率，贪心也不保证整句概率最大，但用它看清流程已经足够。

\subsection{生成循环}

假设每一步恰好选出期望的字符，整个过程如表~\ref{tab:generation} 所示。
\begin{table}[htbp]
\centering
\caption{一次生成轨迹（教学演示）}
\label{tab:generation}
\renewcommand{\arraystretch}{1.2}
\begin{tabular}{clc}
\toprule
步骤 & 右边已经知道的输入前缀 & 新选出的 token\\
\midrule
1 & \bos & 阿\\
2 & \bos\ 阿 & 姆\\
3 & \bos\ 阿姆 & 斯\\
4 & \bos\ 阿姆斯 & 特\\
5 & \bos\ 阿姆斯特 & 朗\\
6 & \bos\ 阿姆斯特朗 & \eos\\
\bottomrule
\end{tabular}
\end{table}

选到结束标记 \eos\ 后停止，去掉特殊标记，显示``阿姆斯特朗''。实际也可设置最大生成长度，避免一直不结束。

\begin{figure}[htbp]
\centering
\begin{tikzpicture}[x=1cm,y=1cm]
\node[flowbox,minimum width=2.4cm,minimum height=1cm] (prefix) at (0,0) {当前输出前缀};
\node[flowbox,minimum width=3cm,minimum height=1cm] (decoder) at (4.2,0) {Decoder\\读取前缀与问题表示};
\node[flowbox,minimum width=2.8cm,minimum height=1cm] (select) at (8.8,0) {Linear + Softmax\\选择下一个 token};
\node[flowbox,minimum width=3cm,minimum height=0.8cm] (m) at (4.2,1.8) {编码一次的问题表示 $\bm M$};
\draw[flowarr] (prefix)--(decoder); \draw[flowarr] (decoder)--(select);
\draw[flowarr] (m)--(decoder);
\draw[flowarr] (select.south)--(8.8,-1.3)--node[below,font=\small]{未结束：追加到前缀，再计算下一步}(0,-1.3)--(prefix.south);
\end{tikzpicture}
\caption{生成循环}
\label{fig:generationloop}
\end{figure}

这解释了``并行计算''和``逐步生成''为什么并不矛盾：对于已知前缀，同一层的各位置可以一起计算；但新的 token 尚未被选出，下一次前缀就还没有确定。实际系统可以缓存已经算过的 Key、Value 来避免部分重复工作，这不改变这里的逻辑。

最后留意一个细节：``阿姆斯特朗''并没有出现在输入问题里，模型照样能输出它——因为最后一层是在整个词表上打分，而不是从输入句子里挑词。至于哪个词该得高分，由训练学到的参数和当前输入共同决定。

\subsection{训练时的输入：Outputs shifted right}

前面的生成过程默认模型已经训练完毕，最后补上这块拼图：训练到底在优化什么。训练时手里已经有问题和正确答案。原图右下角的``Outputs shifted right''，说的就是怎么把正确答案变成 Decoder 的输入：在开头补 \bos，让输入与预测目标错开一位。

\begin{table}[htbp]
\centering
\caption{训练时的输入与目标对齐}
\label{tab:teacherforcing}
\small\renewcommand{\arraystretch}{1.3}
\begin{tabular}{lcccccc}
\toprule
位置 & 1 & 2 & 3 & 4 & 5 & 6\\
\midrule
Decoder 输入 & \bos & 阿 & 姆 & 斯 & 特 & 朗\\
应该预测 & 阿 & 姆 & 斯 & 特 & 朗 & \eos\\
\bottomrule
\end{tabular}
\end{table}

例如输入``姆''的位置，目标是预测``斯''。虽然训练矩阵中右边也放着``斯''，因果掩码阻止这个位置读取它。模型没有偷看当前要预测的答案。

这类使用真实前缀进行训练的方式常称为 teacher forcing。完整答案已知，使不同目标位置的损失能够在一次前向中一起计算；生成时则没有真实前缀可用，只能追加自己已经选出的 token。

\subsection{预测损失}

把源句记为 $s$，目标序列 $y=(y_1,\ldots,y_m)$ 包含结束标记。自回归形式为
\begin{equation}
p_\theta(y\mid s)=\prod_{i=1}^{m}p_\theta(y_i\mid s,y_{<i}).
\label{eq:autoregressive}
\end{equation}
让训练答案获得更高概率，可以等价地最小化负对数似然：
\begin{equation}
\mathcal L(\theta)
=-\frac1m\sum_{i=1}^{m}\log p_\theta(y_i\mid s,y_{<i}).
\label{eq:crossentropy}
\end{equation}
这里给出未使用标签平滑、未考虑 padding 的基本形式；实际批量训练需忽略补齐位置，原论文还使用了标签平滑。

如果某一位置正确 token 的预测概率从 $0.1$ 提高到 $0.8$，这一位置的损失就从 $-\log0.1\approx2.303$ 降到 $-\log0.8\approx0.223$——训练自始至终在做的，就是把正确答案的概率推高。

$\theta$ 包括输入和输出的嵌入、各头的投影、输出混合矩阵、FFN、LN 的缩放和平移、词表输出层等参数。原论文在适用的共享词表设置中还共享了部分嵌入和输出权重；为了把数据流讲清楚，本文没有要求这些符号必须相互独立。

\begin{table}[htbp]
\centering
\caption{可训练参数与当前输入产生的中间结果}
\label{tab:parameters}
\renewcommand{\arraystretch}{1.25}
\begin{tabularx}{\textwidth}{@{}p{2.6cm}Y Y@{}}
\toprule
类别 & 例子 & 它们如何变化\\
\midrule
可训练参数 & 嵌入表、各类 $\bm W$、偏置、$\bm\gamma$ 和 $\bm\beta$ & 根据训练损失更新；正常推理时固定\\
中间结果 & $\bm X,\bm Q,\bm K,\bm V,\bm A,\bm Z,\bm M$ & 每次随当前输入和参数重新计算\\
固定规则 & 正弦位置编码、因果掩码、选定的模型尺寸 & 在本文的设定中不是梯度优化的参数\\
\bottomrule
\end{tabularx}
\end{table}

因此，``训练注意力''并不是直接存下一张永不变化的关联表，而是学习产生投影与权重的参数。最终损失沿着完整计算图反向传播，包括注意力、Add、Norm、FFN 和嵌入。没有人提前告诉每一层的向量应该等于什么数值。

\begin{figure}[htbp]
\centering
\begin{tikzpicture}[x=1cm,y=1cm]
\node[flowbox,minimum width=2.65cm,minimum height=1.15cm] (in) at (0,0) {问题与右移答案\\作为已知输入};
\node[flowbox,minimum width=2.65cm,minimum height=1.15cm] (net) at (4.25,0) {完整 Transformer\\预测各位置概率};
\node[flowbox,minimum width=2.65cm,minimum height=1.15cm] (loss) at (8.5,0) {与正确 token 对照\\计算交叉熵损失};
\draw[flowarr] (in)--(net); \draw[flowarr] (net)--(loss);
\draw[flowarr,dashed] (loss.south)--(8.5,-1.25)--node[below,font=\small]{反向传播：更新模型参数}(4.25,-1.25)--(net.south);
\end{tikzpicture}
\caption{训练的数据流}
\label{fig:training}
\end{figure}

\section{小结}

到这里，Transformer 的每个部件都拆开看过了。串起来再走一遍：它始终在处理一串向量——嵌入把离散编号变成特征，位置编码引入顺序；注意力根据匹配关系汇总不同位置的信息，多头保留多套汇总方式；残差把模块组织成对已有表示的更新，归一化控制尺度，FFN 在每个位置内部进一步构造非线性特征。

左侧将问题编码成 $\bm M$，右侧根据已有答案前缀构造 Query，通过交叉注意力读取 $\bm M$。最后一个位置的向量被映射到词表分数，再变成下一 token 的概率；选出的 token 被追加到前缀，直到结束。训练则用正确答案的概率定义损失，将这整条链上的参数一起调整。

\begin{equation}
\boxed{\text{构造表示}\ \longrightarrow\ \text{交换与变换信息}\
\longrightarrow\ \text{输出条件概率}}
\label{eq:summary}
\end{equation}

最值得记住的不是那张图里有多少个框，而是每个框改变了什么。$\bm A$ 是信息分配比例，不是最终特征；Add 是模型定义的一部分，不是训练完再补上的修正；LN 接收的是前一步计算的结果，不要求括号里永远只有一个字母；FFN 的表达力来自可学习的非线性变换，而不是单纯的升维。

\enlargethispage{20pt}

``Transformer''也并不强制所有模型都必须保留原图中的两座塔。这里讨论的是原始编码器—解码器版本；理解它的模块和数据流以后，再看使用其中部分结构的模型，就可以逐项比较它保留了哪些连接、修改了哪些约束，而不必重新从术语开始。

\end{document}
